% =====================================================================
% APPENDIX A (MLLB, Maria Martinez): use of AI code assistants
% Sources: docs/AI_USAGE_LOG.md (Bounty candidates B6, B2/B3/B7/B11, B10,
% B19; one report per error), CHATS protocol.
% =====================================================================
\appendixbanner{APPENDIX A: USE OF AI CODE ASSISTANTS}
\addcontentsline{toc}{section}{Appendix A: Use of AI code assistants}\label{app:ai-code}

\subsection*{Workflow}

For the code, each block went through the same cycle with Claude: a written proposal of the decisions,
which we accepted or changed; code and tests written by the assistant and run by it on a CPU copy with
synthetic data; then our own run of the tests, the smoke test and the real runs on the laptop, whose output
we pasted back; and finally the analysis, in which the assistant recomputed the metrics from the saved
predictions and we checked the figures and tables. For my blocks (data, evaluation, recurrent networks,
self-attention and Transformer) I read each function with the assistant's explanation before running it.
Every session is in \texttt{docs/AI\_USAGE\_LOG.md}.

\subsection*{Errors detected (candidates for the AI Error Bounty)}

Each assistant error of the project is documented in exactly one of our four reports. These four sheets
cover the code errors of my blocks, of the shared zero-shot experiment and of the project's file handling:
\begin{itemize}[leftmargin=1.2em,itemsep=2pt]
  \item \textbf{A code comment that contradicted the code} in my Transformer model (B6; Sheet C1).
  \item \textbf{Files reported as written when an older copy had reached the folder}, four times in four
        blocks (B2, B3, B7, B11; Sheet C2).
  \item \textbf{Tensors on two devices} in the zero-shot code (B10; detected by Oriana and me; Sheet C3).
  \item \textbf{A clean compilation reported without reading the BibTeX log} (B19; Sheet C4).
\end{itemize}
The errors found in Oriana's blocks and runs (out of memory when generating, a figure legend over a data
point, commands without the repository folder, a wrong cause for a log flag, a smoke test that stopped)
are documented in her MLLB report.

\subsection*{AI Error Bounty sheets}

\subsubsection*{Sheet C1. A comment that contradicted the code (Claude)}
\begin{description}[style=nextline,leftmargin=1.2em]
  \item[Prompt] The complete self-attention and Transformer block (decision D18): model, experiments and
        tests.
  \item[Faulty response] The comment above \texttt{register\_buffer("pe", \ldots, persistent=False)} in
        \file{src/models/transformer.py} said that the positional-encoding buffer is ``saved with the model''.
  \item[Detected by] The assistant, when it checked the PyTorch details of my block note against the code
        after the runs.
  \item[Proof] With \texttt{persistent=False} PyTorch leaves the buffer out of \texttt{state\_dict()}, so it
        is not in the checkpoint: \texttt{'pe' in TransformerClassifier().state\_dict()} is \texttt{False},
        while the attribute \texttt{pe} exists.
  \item[Correction] Comment rewritten: the buffer moves to the GPU with the model, is not trained, and is
        not saved, because it is rebuilt when the model is created. The behaviour of the code did not change,
        so the runs and the results stand.
  \item[Reproduction] \texttt{python -c "from src.models.transformer import TransformerClassifier as T;
        print('pe' in T().state\_dict())"} prints \texttt{False}.
  \item[Reflection] A wrong comment is dangerous in a course with an oral defence about specific lines: it
        would have taught me the wrong behaviour. Buffers, parameters and the state dictionary are now one of
        the PyTorch points I can explain with this example.
\end{description}

\subsubsection*{Sheet C2. Older copies of files left in the folder (Claude)}
\begin{description}[style=nextline,leftmargin=1.2em]
  \item[Prompt] Four times, in four blocks, to write files to the project folder after an edit: the EDA
        decision D12 (B2), the evaluation decisions D14 in my evaluation block (B3), three files of the
        generative block (B7), and \file{src/zero_shot.py} with \file{docs/DECISIONS.md} (B11).
  \item[Faulty response] Each copy was reported as successful, or was about to be used, while the folder
        still held the previous content, because the file was copied right after being written, before the
        new version could be read.
  \item[Detected by] The assistant itself each time, when it re-read the folder before the next step; from
        B7 on, through the rule we added to the common protocol (check every written file in the folder).
  \item[Proof] Sizes in the folder that did not match the intended files (for D14: 22,839 bytes instead of
        28,866, with the assistant's own modification time, so no other conversation had changed it).
  \item[Correction] Each file copied again after waiting, and compared by size or byte by byte with the
        intended version before anyone was told it was written or asked to run anything.
  \item[Reproduction] Not in git, because no incomplete copy was ever committed; the evidence is the sizes
        recorded in \texttt{docs/AI\_USAGE\_LOG.md} (B2, B3, B7, B11). The check is repeated after every write.
  \item[Reflection] The same error came back four times until checking became a rule rather than an
        intention: an assistant that writes files must prove that the right version is there.
\end{description}

\subsubsection*{Sheet C3. Tensors on two devices (Claude)}
\begin{description}[style=nextline,leftmargin=1.2em]
  \item[Prompt] The zero-shot experiment with CLIP on CIFAR-100 and on ECG images (shared task T24).
  \item[Faulty response] The image features were moved to the CPU while the text features stayed on the
        GPU, and the two were multiplied. The assistant's tests used a fake CLIP on the CPU only, so they
        could not show it.
  \item[Detected by] Oriana and me, when the run stopped on the laptop.
  \item[Proof] On the laptop the run stopped with \texttt{RuntimeError: Expected all tensors to be on the
        same device}.
  \item[Correction] Both feature sets are returned on the CPU, with a regression test. In the same fix the
        model was loaded as \texttt{ViT-B-32-quickgelu}, the exact activation of the OpenAI weights used in the
        class notebook.
  \item[Reproduction] Remove \texttt{.cpu()} from the text features and from the similarity function of
        \file{src/zero_shot.py} and run \texttt{python -m src.zero\_shot run --fast} on the GPU: same error.
  \item[Reflection] Tests on the CPU cannot catch device errors; the smoke run on the real GPU is a test in
        its own right and must not be skipped.
\end{description}

\subsubsection*{Sheet C4. A clean compilation reported without the BibTeX log (Claude)}
\begin{description}[style=nextline,leftmargin=1.2em]
  \item[Prompt] Compile the reports and check that they have no errors (the condition for closing a report).
  \item[Faulty response] Oriana's two reports were reported three times as compiling with ``0 errors, 0
        undefined'', checking only the LaTeX log. The BibTeX log had one error: an entry of
        \file{report/references.bib} contained a comment, which BibTeX does not accept inside an entry, so the
        rest of that entry was skipped. The reference still printed, so the PDF looked correct.
  \item[Detected by] The assistant, when it compiled my reports and read the \texttt{.blg} file as well.
  \item[Proof] The BibTeX logs of the earlier builds ended with ``(There was 1 error message)'' and pointed to
        the line of the entry.
  \item[Correction] The comment was moved out of the entry; every compilation since then checks the LaTeX log
        and the BibTeX log (0 errors in the four reports).
  \item[Reproduction] Put a \texttt{\% comment} line inside the entry \texttt{chelba2017ngram}, run BibTeX on
        any report and read the end of its \texttt{.blg}.
  \item[Reflection] ``It compiles'' has two logs; a check is only as good as the file it reads.
\end{description}

\subsection*{Reflection}

The coding assistant was most useful for long code that can be checked: experiment scripts, tests,
summaries, tables and figures. Its errors fell where its own checks could not reach: the real GPU, the
real folder, and statements it did not verify. The rules of Section~\ref{sec:ai-env}, above all running
everything ourselves and generating every number from the logs, are what kept those errors out of the
results.
